\section{Context Problem：现实工作不是一段完整 prompt}

\videofigure{figures/context-experiment.jpg}{GDPval 主动缩短 prompt、移除部分上下文后，模型胜率明显下降，暴露 context acquisition 的重要性。}{00:45:48--00:46:48}

\videofigure{figures/context-problem.jpg}{仓库维护者、领域专家与外部承包者之间可有数倍效率差异，因为真实工作依赖隐性上下文。}{00:46:48--00:47:32}

真实员工知道历史决策、组织偏好、文件位置、例外流程和“谁能回答什么”。Benchmark 常把这些信息预先整理进 prompt，等于跳过了 agent 最困难的工作之一：发现缺什么、向谁询问、到哪里找、怎样判断信息仍然有效。

\begin{warningbox}{上下文不足与能力不足容易混淆}
一个聪明但没有仓库、客户和组织背景的 agent，可能输给能力普通但上下文充分的人。评估应区分 reasoning、context retrieval、权限和环境接入造成的失败。
\end{warningbox}

\begin{importantbox}{下一代 agent 要学习 context acquisition}
Agent 不应被动等待完美 prompt，而要主动列出缺失信息、检索内部资料、询问澄清问题并维护来源与时效。这是从 benchmark solver 走向真实同事的关键能力。
\end{importantbox}

\videofigure{figures/gdpval-implications.jpg}{GDPval 更支持近期的专业辅助与人类监督，而不是无条件的全自动替代。}{00:47:32--00:49:10}

\subsection{本章小结}

真实工作的瓶颈常在获取上下文，而非生成最终文本。部署与评估都应把检索、权限、澄清和组织记忆纳入 agent loop。

